\chapter{第二十三讲\quad 贪污贿赂罪}
\label{lecture:23}
\noindent{\small\hyperref[ref:lecture:23]{查阅本讲扩展规则}}\par

本讲把一笔钱分成「原来是谁的、经谁的权力、最后归谁」三问。记 \(a\) 为收财或挪款者，\(b\) 为请托人，\(c\) 为最终实施职务行为者。\(Q_i\) 总指具体罪名及行为时点的数额或情节条件。法释〔2026〕6号自2026年5月1日起施行，其中与此前解释不一致的事项按该解释核对，跨时点案件先处理适用法源问题。

\section{第一节\quad 贪污罪}

\paragraph{法条与模型}
令 \(N(a)\) 为国家工作人员，\(V(a,w)\) 为利用职务便利接触或控制财物 \(w\)，\(P(w)\) 为公共财物，\(J(a,w)\) 为非法占有目的，\(E(a,w)\) 为侵吞、窃取、骗取或以其他手段非法占有。第382条第一款的模型为
\[
N(a)\land V(a,w)\land P(w)\land J(a,w)\land E(a,w)
\Longrightarrow B_{382,1}.
\]
第383条再核对处罚所需数额、情节。第382条第二款另有受委托管理、经营国有财产的主体分支，不能用公共财物概念无条件替换国有财物。第三款处理与前两款人员勾结、伙同贪污。

\paragraph{推导}
受委派在国有企业从事公务的会计甲，代表单位收回客户货款，随后用虚假对账将该笔款留作己用。由委派和公务事实确定 \(N(a)\)，代收职责给出 \(V\)，已代表单位收取的款项为 \(P\)，隐匿和据为己有给出 \(J,E\)。依第382条第一款得构成要件结论，再依第383条及司法解释核对 \(Q_{383}\) 与档次。

\paragraph{边界}
甲只是凭工作地点熟悉程度撬柜拿钱，且取财与其职责权限无关，不能用「在单位工作」推出 \(V\)，须另查盗窃等罪。若以集体研究名义仅向领导和管理层秘密私分国有资产，2026年解释（二）第18条第一款规定按贪污评价。不能单凭「集体决定」推出第396条私分国有资产罪。

\section{第二节\quad 挪用公款罪}

\paragraph{法条、立法解释与定义}
令 \(F\) 为公款，\(U\) 为归个人使用，\(I\) 为进行非法活动，\(Y\) 为进行营利活动，\(t\) 为未还时长，\(Q\) 为相应数额条件。第384条的三个分支写为
\[
N\land V\land F\land U\land
\bigl[(I\land Q_I)\lor(Y\land Q_Y)\lor
(Q_T\land(t>3\text{个月})\land\mathrm{NotReturned})\bigr]
\Longrightarrow B_{384}.
\]
\(Q_I,Q_Y,Q_T\) 分别按各自的司法解释入罪标准处理。全国人大常委会对第384条第一款的解释列举三个 \(U\) 的入口。公款供本人、亲友或其他自然人使用，或者以个人名义供其他单位使用，或者个人决定以单位名义供其他单位使用且谋取个人利益。2026年解释（二）第9条将个人以虚构付款事由、账款不按规定入账等逃避监管的方式将公款供其他单位使用，纳入「以个人名义」。

\paragraph{推导}
主管公款的国家工作人员甲私自把公款借给朋友经营商店，金额已达营利型标准。身份、主管与资金属性给出 \(N,V,F\)，供其他自然人使用给出 \(U\)，经营商店给出 \(Y\)，核定金额给出 \(Q_Y\)。第二个析取项已经为真，故即使两个月归还，仍不能以第三分支的三个月条件否定该案。

\paragraph{边界}
「未归还」与「一开始意欲占有」是不同事实。经营失败后客观不能退还，不能只据失败反推贪污故意。若查明行为人从一开始便伪造销账并欲永久排除单位控制，则转审第382条。数额巨大不退还按2026年解释（二）第10条核对公诉前、客观原因和职权追回情形。

\section{第三节\quad 受贿罪}

\paragraph{法条与模型}
\(V_{a,c}\) 表示利用本职权，或者通过对实施者 \(c\) 存在职务隶属、制约关系而利用职务便利。令 \(S\) 为索财，\(R\) 为非法收财，\(M\) 为为他人谋取利益，\(L\) 为收财与职务行为的联系。第385条第一款可按索取、收受两支记忆。
\[
N(a)\land V_{a,c}\land L\land
\bigl[S\lor (R\land M)\bigr]\Longrightarrow B_{385}.
\]
此式抽出基本构造，另核数额和情节。第二款还规定经济往来中违反国家规定，利用职务便利收受回扣、手续费归个人所有的分支。\(M\) 的证明可来自实际谋利、承诺、明知具体请托事项等法定解释路径。履职时未被请托，事后基于该履职事由收财，按2016年贪污贿赂解释第13条核对。

\paragraph{推导}
负责招标审批的甲收了请托人乙交付的款，并答应利用自己的审批权加速审查。甲的身份和审批权给出 \(N,V\)，收款给出 \(R\)，答应办事给出 \(M,L\)。依第385条和上述解释推出收受型候选，已证明实际收受且数额条件满足时，不以加速审批最终完成为必要。若是甲通过有职务隶属、制约关系的其他国家工作人员办事，依2026年解释（二）第14条仍可认定职务便利。

\paragraph{计价与边界}
收受财物一般按收受时价值 \(v(t_{\rm in})\) 认定。以股票、股权预期收益作为贿赂形式并构成犯罪的，按2026年解释（二）第11条区分案发时已获利和尚未获利，前者按实际获利，后者一般按案发时市场价格与支付价格的溢价。普通朋友送礼且证实没有职务联系，不能由价格倒推 \(L\)。受贿又构成相关渎职罪时，通常按2016年解释第17条审查并罚，有刑法特别规定时优先用特别规则。

\section{第四节\quad 行贿罪}

\paragraph{法条与模型}
设 \(U_b\) 为谋取不正当利益，\(G_{b,a}\) 为给予国家工作人员财物，\(D\) 为被勒索给予财物，\(A\) 为获得不正当利益。第389条第一款的一般路径为 \(U_b\land G_{b,a}\)，另核第390条与解释的数额和情节。第三款设 \(D\land\neg A\) 的不构成行贿规则。第二款另查经济往来中违反国家规定给财物或回扣、手续费的情形。单位行贿、第391条对单位行贿、第390条之一对有影响力的人行贿分别另立主体与相对人条件。

\paragraph{推导}
乙送给负责审批的国家工作人员甲一笔达标款，请甲帮助绕过法定准入条件。给付事实得 \(G_{b,a}\)，绕过法定条件的目的得 \(U_b\)，依第389条推出乙的行贿候选。甲是否构成受贿还要分别证明其职务联系及故意，不能把乙的 \(U_b\) 直接复制为甲的事实。

\paragraph{边界}
将乙改为在甲勒索下付款，仅领取本应取得的法定证照，证据证明 \(D\land\neg A\)，依第389条第三款不成立行贿。单位主管决定付款的，仍须查利益实际归属。依2026年解释（二）第16条，决定主体和利益归单位是单位行贿的重要判据，不能只看使用哪一个账户。刑法修正案（十二）调整第390条，复习从重以及被追诉前主动交代的特殊从宽时，应逐项查现行条文。

\section{第五节\quad 斡旋受贿}

\paragraph{法条与路径图}
令 \(C_{a,c}\) 表示甲对丙有职务隶属、制约关系，\(F_{a,c}\) 表示利用本人职权、地位形成的便利以通过丙的职务行为办事。第388条的典型路径可写为
\[
b\xrightarrow{\text{请托、财物}}a
 \xrightarrow{F_{a,c}}c
 \xrightarrow{\text{职务行为}}b .
\]
\(a,c\) 均为国家工作人员，\(a\) 利用该便利，通过 \(c\) 的职务行为为 \(b\) 谋取不正当利益并索取或收受其财物，才进入该条分支。若实质是 \(C_{a,c}\) 所体现的本职便利，则优先按第385条对应事实评价。2026年解释（二）第13、14条分别细化承诺和职务制约的认定。

\paragraph{推导}
甲以其任职地位形成的工作联系，承诺通过无职务制约关系的另一单位国家工作人员丙，为乙取得不正当许可，并收下乙的财物。任职联系证实 \(F_{a,c}\)，请托事项证实不正当利益，承诺、收财和职务路径均已查清。在数额或情节达到相应标准且其余条件成立时，依解释（二）第13条和第388条，以受贿论处。甲尚未向丙转达，请托事项未转达这一事实不足以否定本案结论。

\paragraph{边界}
将丙换成实际上受甲职务制约的国家工作人员，适用入口要重新分流。教材用上下级与同级帮助记忆，但「同级」或「非直接上下级」本身既不能证明 \(F_{a,c}\)，也不能否定 \(C_{a,c}\)。应核对任职单位、权责安排及实际制约。丙若对交易不知情，不能从其办事推出共同受贿。

\section{第六节\quad 利用影响力受贿罪}

\paragraph{法条与模型}
第388条之一关注特定影响关系 \(I\)、法定职务路径 \(V\)、不正当利益 \(U\)、索取或收受财物 \(G\)、数额较大或其他较重情节 \(Q\)。近亲属或其他关系密切者可通过该国家工作人员的职务行为，也可利用该国家工作人员职权或地位形成的便利，经其他国家工作人员办事。离职国家工作人员及其近亲属、关系密切者另核对原有职权或地位形成的便利。模型为
\[
I\land V\land U\land G\land Q\Longrightarrow B_{388\text{之一}} .
\]

\paragraph{推导}
某负责审批的国家工作人员丙的近亲属甲，确实借助与丙的关系，促成丙给请托人乙违法审批。甲以此私自收取乙的财物，丙不知道甲收钱，已达法定数额。近亲属和审批路径得 \(I,V\)，违法审批得 \(U\)，收财、数额得 \(G,Q\)，推出利用影响力受贿。无意思联络，不能认定甲与丙共同收受。若甲在介绍贿赂中占款，依2026年解释（二）第17条还要分别检查该条规定的截留、撮合或共同贿赂分支。

\paragraph{边界}
改成甲虚构与丙关系密切，以非法占有为目的骗取乙财物，真实影响关系 \(I\) 未成立。达到诈骗构成要件时，依解释（二）第17条按诈骗评价。若仅撮合贿赂双方，未独立利用影响力收财，则转查第392条及严重情节。共同实施行贿或受贿同时构成介绍贿赂罪的，核对该解释的择重规则。

\paragraph{普通罪名：巨额财产来源不明}
依第395条第一款，先由国家工作人员身份、查明的财产与支出、合法收入等证据核定差额 \(D\)，再查是否达到差额巨大及能否说明来源。不得把同一笔资金形成的支出和财产余额重复计入差额。2026年解释（二）第5条规定，三百万元以上、不满一千万元是「差额巨大」一档。设案件中经核对且避免重复的差额为四百万元，依法责令说明后不能说明来源，且其他要件均具备，则可依第395条第一款的该档评价。若核实其中两百万元为合法来源，剩余差额为两百万元，原金额分支随即失效。

\paragraph{闭卷记忆}
\begin{enumerate}
\item 先画钱的箭头。公共财物出单位后，问行为人意欲占有还是暂时使用。暂用再分非法、营利、超期未还。
\item 再画办事的箭头。自身权力或实质制约看普通受贿，职权地位形成便利经他人办事看斡旋，近亲或原有地位等特定影响看第388条之一。
\item 最后对给财方单独画图，依次核对对象、获利目的、是否被勒索、实际利益归属，以及阈值和例外。
\end{enumerate}
口诀「\textbf{公款问占用，交易问权路。收给分开推，数额末尾核}」。用「本职、地位、关系」分别回忆 \(V_{a,c},F_{a,c},I\)。

\paragraph{闭卷自测与答案}
\begin{enumerate}
\item 挪款给朋友营利，数额达标、两个月归还，能以不足三个月否定挪用公款罪吗？答：不能。营利、数额已使第二个析取项成立，第三个析取项的时长不支配它。
\item 甲有权制约其他单位国家工作人员乙，甲收款并通过乙办事，是否一定是斡旋受贿？答：先核实是否为职务上的隶属或制约。若是且其余要件具备，属第385条利用职务便利；不能按单位不同直接选择第388条。
\item 熟人确实与公职人员关系密切，与仅假冒这种关系相比，哪个谓词不同？答：真实影响关系 \(I\)。两案还须独立检查职务路径、收财、非法占有目的、欺骗及数额等要件。
\end{enumerate}
